% GraphHunter --- arquitectura (preambulo en español).
%
% Hereda la estructura del preambulo ingles pero localiza las etiquetas
% de environments (teorema/lema/demostracion/...) y los kw de algorithm2e
% (para/mientras/si/entonces/...). Macros propias (`\Ord`, `\filepath`,
% etc.) se mantienen con el mismo nombre.

\usepackage[utf8]{inputenc}
\usepackage[T1]{fontenc}
\usepackage[spanish,es-nodecimaldot]{babel}
\usepackage{lmodern}

\usepackage[a4paper,margin=2.4cm]{geometry}

\usepackage{amsmath, amssymb, amsthm}
\usepackage{mathtools}

% Pseudocodigo.
\usepackage[ruled,linesnumbered,algochapter]{algorithm2e}
\SetKwInput{KwIn}{Entrada}
\SetKwInput{KwOut}{Salida}
\SetKwComment{Comment}{$\triangleright$ }{}
\SetKwFor{For}{para}{hacer}{fin}
\SetKwFor{While}{mientras}{hacer}{fin}
\SetKwFor{ForEach}{para cada}{hacer}{fin}
\SetKwIF{If}{ElseIf}{Else}{si}{entonces}{sino si}{sino}{fin}
\SetKwRepeat{Repeat}{repetir}{hasta}
\SetKwSwitch{Switch}{Case}{Other}{segun}{hacer}{caso}{en otro caso}{fin caso}{fin segun}
\SetKw{KwTo}{a}
\SetKw{KwBreak}{romper}
\SetKw{KwContinue}{continuar}
\SetKw{KwTrue}{verdadero}
\SetKw{KwFalse}{falso}
\newcommand{\Break}{\KwBreak}
\newcommand{\Continue}{\KwContinue}
\renewcommand{\Return}{\KwRet}  % mantiene "\Return" pero imprime "devolver"
\SetKw{KwRet}{devolver}

% Listings para snippets Rust (iguales al EN).
\usepackage{listings}
\usepackage{xcolor}
\definecolor{rustcomment}{rgb}{0.35,0.45,0.45}
\definecolor{rustkw}{rgb}{0.55,0.10,0.55}
\definecolor{ruststr}{rgb}{0.10,0.45,0.10}
\lstdefinelanguage{Rust}{
    morekeywords={fn,let,mut,pub,struct,enum,impl,trait,where,if,else,for,while,
                  loop,match,return,use,mod,const,static,self,Self,as,in,ref,move,
                  async,await,dyn,unsafe,crate,super,extern,type,break,continue},
    morecomment=[l]{//},
    morecomment=[s]{/*}{*/},
    morestring=[b]",
    morestring=[b]',
}
\lstset{
    basicstyle=\ttfamily\small,
    keywordstyle=\color{rustkw}\bfseries,
    commentstyle=\color{rustcomment}\itshape,
    stringstyle=\color{ruststr},
    showstringspaces=false,
    frame=single,
    numbers=none,
    captionpos=b,
    breaklines=true,
    tabsize=4,
}
\lstdefinestyle{rust}{language=Rust}

% Hyperref + bookmarks.
\usepackage{hyperref}
\hypersetup{
    colorlinks=true,
    linkcolor=blue!50!black,
    citecolor=black,
    urlcolor=blue!50!black,
    pdftitle={GraphHunter: arquitectura (referencia estilo CLRS)},
    pdfauthor={Lucas Sotomayor / Base4Security},
    bookmarks=true,
    bookmarksopen=true,
    bookmarksnumbered=true,
}
\usepackage{bookmark}
\usepackage{booktabs}
\usepackage{tikz}
\usetikzlibrary{arrows.meta, positioning, shapes.geometric, fit, backgrounds}

% Entornos teorema (etiquetas en español).
\theoremstyle{plain}
\newtheorem{theorem}{Teorema}[chapter]
\newtheorem{lemma}[theorem]{Lema}
\newtheorem{corollary}[theorem]{Corolario}

\theoremstyle{definition}
\newtheorem{definition}[theorem]{Definición}

\theoremstyle{remark}
\newtheorem*{invariant}{Invariante}
\newtheorem*{note}{Nota}

% "Demostración" en vez de "Proof".
\renewcommand{\proofname}{\textit{Demostración}}

\newenvironment{complexity}{%
    \par\vspace{0.4em}%
    \noindent\textbf{Complejidad.}\nopagebreak\par\nopagebreak\vspace{0.2em}%
}{\par\vspace{0.4em}}

\newenvironment{inthecode}{%
    \par\vspace{0.4em}%
    \noindent\textbf{En el código.}\nopagebreak\par\nopagebreak\vspace{0.2em}%
}{\par\vspace{0.4em}}

% Macros (mismas que EN).
\newcommand{\Ord}[1]{\ensuremath{O(#1)}}
\newcommand{\Thet}[1]{\ensuremath{\Theta(#1)}}
\newcommand{\filepath}[1]{\texttt{\small\detokenize{#1}}}
\newcommand{\code}[1]{\texttt{\detokenize{#1}}}
\newcommand{\StrId}{\ensuremath{\mathsf{StrId}}}
\newcommand{\Compact}{\ensuremath{\mathsf{CompactRelation}}}
\newcommand{\Graph}{\ensuremath{\mathsf{InMemoryGraph}}}
\newcommand{\Scorer}{\ensuremath{\mathsf{AnomalyScorer}}}

\makeatletter
\renewcommand{\@algocf@pre@ruled}{\hrule height.8pt depth0pt \kern2pt}
\renewcommand{\@algocf@post@ruled}{\kern2pt\hrule height.8pt\relax}
\makeatother

\usepackage{titlesec}
\titleformat{\chapter}[display]
  {\normalfont\huge\bfseries}
  {\chaptertitlename~\thechapter}
  {20pt}
  {\Huge}

\newcommand{\passthrough}[1]{#1}
